% answers_short.tex —— 2006 年·全国卷Ⅰ·文科 · 简短答案速查
% 由 tools/make_short_answers.py 自 answers.tex 的 \gksolhead/\gkanswer 提取，
% 并与 2006/2006_数学真题_全国1文_答案要点.md 逐题比对后生成，勿手改（改源 MD 后重跑）。

\begin{gknotice}
\textbf{使用说明}：本节只给最终答案便于快速核对，完整推导见解析版。答案为编写者独立推导，\textbf{非官方标准答案}。
\end{gknotice}

\gkpart{选择题}
\begin{center}\begin{tabular}{ccccccccccccc}
\toprule
题号 & 1 & 2 & 3 & 4 & 5 & 6 & 7 & 8 & 9 & 10 & 11 & 12 \\
\midrule
答案 & C & B & D & A & D & C & B & C & C & A & A & B \\
\bottomrule
\end{tabular}\end{center}

\gkpart{填空题}
\begin{itemize}[leftmargin=2.6em]
  \item[\textbf{13.}] $\dfrac12$
  \item[\textbf{14.}] $\dfrac\pi3$
  \item[\textbf{15.}] $11$
  \item[\textbf{16.}] $2400$
\end{itemize}

\gkpart{解答题}
\begin{itemize}[leftmargin=2.6em]
  \item[\textbf{17.}] $a_n=2\cdot3^{n-3}$ 或 $a_n=2\cdot3^{3-n}$。
  \item[\textbf{18.}] $A=\dfrac\pi3$ 时取得最大值，最大值为 $\dfrac32$。
  \item[\textbf{19.}] (1) $\dfrac49$；(2) $\dfrac{604}{729}$
  \item[\textbf{20.}] (1) 见解析版；(2) $\sqrt{\dfrac23}$
  \item[\textbf{21.}] $2$（$1<a<\sqrt2$）；$\dfrac{a^2}{\sqrt{a^2-1}}$（$a\geqslant\sqrt2$）。
  \item[\textbf{22.}] $a\in\left(-\infty,-\sqrt{\dfrac32}\right]\cup[1,+\infty)$。
\end{itemize}
